% ── Rakouský model: Kammer-system a Allgemeinverbindlicherklärung ────────────
je pro \acacc{ČR} nejbližší praktickou inspirací, protože poukazuje na rozdíl mezi odborovou organizovaností a~pokrytím \aclins{KS}.
Rakouský model je~založen na~silném odvětvovém vyjednávání a~vysoké organizovanosti zaměstnanců i~zaměstnavatelů díky povinnému
členství v~hospodářské komoře (\textit{Wirtschaftskammer Österreich, WKÖ}) a~komoře pracovníků (\textit{Arbeiterkammer, AK}).
Členství v~odborech (největší sdružení \textit{Österreichischer Gewerkschaftsbund, ÖGB}) není povinné.\par
V~praxi to znamená, že pokrytí \acs{KS} zůstává dlouhodobě velmi vysoké i~při nižší odborové
organizovanosti, než je běžné v~severských zemích. Podniková úroveň zde typicky
plní často implementační roli vůči odvětvovým standardům.\par

V~datech \acs{OECD}~\acs{ICTWSS} se~Rakousko dlouhodobě drží u~hodnot pokrytí
blízkých \SI{98}{\percent} (vyšší v~soukromé sféře), zatímco
odborová organizovanost je výrazně nižší (přibližně \SI{26}{\percent}).
Tento případ ukazuje, že vysokého pokrytí \acs{KS} lze dosáhnout s nižší odborovou organizovaností.\par

Vysoké pokrytí zároveň snižuje prostor pro~vnitrostátní mzdový dumping mezi
podniky téhož odvětví, protože minimální mzdy jsou stanovovány odvětvovými smlouvami.
To je rozdíl vůči \acloc{ČR}, kde je odvětvová vrstva velmi slabá a~celkové pokrytí \acs{KS} nízké.
Minimální mzda stanovená zákonem pro všechny sektory hospodářství neodráží odvětvové rozdíly, které se prohlubují.\par

Převzetí rakouského komorového systému je v~\acloc{ČR} politicky i~ústavněprávně problematické. Přenositelná je však logika vyššího pokrytí prostřednictvím odvětvového vyjednávání a~funkční procedury rozšíření závaznosti -- pokud se má zvýšit pokrytí \aclgen{KV}, nelze spoléhat jen na zakládání nových odborových organizací v~jednotlivých podnicích a~organizaci zaměstnavatelů ve svazích. Český \S~7 zákona o~kolektivním vyjednávání právní základ poskytuje, ale jeho praktické využití zůstává omezené. Přenositelná je proto zejména procesní dostupnost, informační opora a~politická vůle nástroj používat.~\cite{Eurofound_AT_Actors2024}~\cite{Eurofound_AT_CB2024}\par